\documentclass[a4paper, 12pt]{article}

\usepackage[frenchb]{babel} % Document en français
\usepackage[utf8]{inputenc} % Document au format UTF8
\usepackage[T1]{fontenc} % Suppression d'un warning pour la langue francaise
\usepackage[left=2.2cm, right=2.2cm, top=2.5cm, bottom=2.5cm]{geometry} % Mise en page
\usepackage{setspace} % Permet de définir l'interligne
\usepackage{titling}
\usepackage{graphicx}

\singlespacing % Interligne à 1

\newcommand{\HRule}{\rule{\linewidth}{0.5mm}}
\graphicspath{{res/}}

\usepackage{newtxtext,newtxmath} % pour la police time


\begin{document}

\begin{titlepage}
\begin{center}

% Upper part of the page. The '~' is needed because \\
% only works if a paragraph has started.
\includegraphics[width=0.15\textwidth]{./sedes_ucl.png}~\\[1cm]

\textsc{\LARGE Université Catholique de Louvain}\\[1.5cm]

\textsc{\Large Travail pratique de Management humain}\\[0.5cm]

% Title
\HRule \\[0.4cm]
{ \huge \bfseries Les rôles de la fonction Ressources Humaines\\[0.4cm] }

\HRule \\[1.5cm]

% Author and supervisor
\noindent
\begin{minipage}[t]{0.4\textwidth}
\begin{flushleft} \large
\emph{Auteurs:}\\
Ivan \textsc{Ahad} \\
Aymeric \textsc{De Cocq} \\
Solaiman \textsc{El Jilali} \\
Léonard \textsc{Julémont} \\
Eddy \textsc{Ndizera} \\
Florian \textsc{Thuin}

\end{flushleft}
\end{minipage}%
\begin{minipage}[t]{0.4\textwidth}
\begin{flushright} \large
\emph{Professeurs :} \\
Nathalie \textsc{Delobbe} \\
Emilie \textsc{Malcourant}
\end{flushright}
\end{minipage}

\vfill

% Bottom of the page
{\large \today}

\end{center}
\end{titlepage}

%\maketitle

%\tableofcontents

\section{Analyse des rôles selon les discours}
Notre groupe a eu l'occasion de participer à la table ronde accueillant trois directeurs de ressources humaines : Christine Thiran, Frédéric Nolf, Alain Goergen.

Christine Thiran travaille actuellement comme directrice des ressources humaines pour le groupe Mestdagh. Cela fait cinq mois qu'elle est à la tâche dans cette entreprise familiale du secteur de la grande distribution. La grande distribution est un secteur en grande difficulté disposant de marges limitées qui emploie majoritairement du personnel non-qualifié. Le groupe Mestdagh a comme vision d'offrir de l'emploi dans des zones économiquement défavorisées, et doit actuellement intégrer 16 anciennes enseignes (et leurs employés) du groupe Carrefour.

Elle a pour rôle de développer une nouvelle culture d'entreprise qui pourrait convenir au personnel du Groupe Mestdagh et au personnel des magasins Carrefour à intégrer. Ce rôle s'articule autour de plusieurs missions : être l'initiatrice/le moteur du changement, garder l'esprit familial de l'entreprise et sa proximité avec les gens, offrir des perspectives de carrière au sein des différents magasins. Son défi principal est très lié à son rôle puisque c'est de trouver un moyen de mettre d'accord les points de vues divergents et les différentes pratiques des employés des magasins Carrefour avec celles des magasins Mestdagh. Les pratiques mises en place pour y arriver sont entre autres : l'évaluation des compétences, l'écoute des ressentis des travailleurs et de leurs propositions, la concertation syndicale, la coordination avec le CEO.

Frédéric Nolf est responsable des ressources humaines pour IBA depuis 2007. IBA est une société côtée en bourse, basée sur l'innovation. Lors de sa création c'était une spin-off de l'UCL et elle est actuellement encore en alliance avec les universités pour trouver des nouveaux talents. Elle est côtée en bourse. Son CEO a changé il y a trois ans.

Son rôle est d'assurer à l'entreprise qu'elle dispose des bonnes personnes pour son développement et que les personnes soient dans la bonne entreprise pour leur développement. Ce rôle s'articule autour de plusieurs missions : s'assurer de la motivation de son personnel, assurer un \og{} sens du focus\fg{}(nous y reviendrons), créer des liens de compétence à travers les divers départements d'IBA.


Dans le cadre de ces missions, il doit relever de nombreux défis : faire rêver les ingénieurs nouvellement diplomés de travailler chez IBA, s'assurer que ces ingénieurs s'intègrent à la culture IBA, s'assurer que les ingénieurs travaillent réellement sur ce qui leur est demandé (\og{} quand il y a beaucoup de gens très intelligents, chacun a son idée sur ce qu'il faut faire \fg{}), informer les membres du personnel à travers le monde sur les projets actuels pour trouver les membres les plus aptes à les résoudre.

Ces défis nécessitent de mettre en place certaines pratiques de GRH : faire des enquêtes de satisfaction et y répondre de manière concrète, mettre en place des activités communes (sport), prendre en compte toutes les idées proposées, créer des projets qui lient les convictions personnelles du personnel à la culture de la société (mise en place d'une cellule de réflexion pour le développement durable).

Alain Goergen est directeur des ressources humaines pour la police fédérale.

% TODO : Analyse de Alain Goergen

\section{Comparaison des rôles RH}
\section{Comparaison des rôles RH}

Maintenant que nous avons établi les rôles des différents intervenants par rapport à leur discours, nous pouvons engager une comparaison de ceux-ci. \newline

%Comparaison

Pour ce qui est d'Alain Goergen, on voit directement son rôle de gestionnaire du changement à la Police Fédérale. Il s'agit ici d'un rapprochement que l'on peut faire avec Christine Thiran qui doit aussi gérer le changement qui s'opère actuellement chez Mestdagh. Néanmoins, ces deux types de changements sont différents à la Police et chez Mestdagh. M. Goergen a du restructurer le département RH et par la suite les autres départements, ainsi que gérer les inquiétudes de chacun quant à ce changement. Ce changement se fait cependant en incluant les mêmes personnes qui ont déjà adopté la culture organisationnelle de la police. Par contre, dans le cas des magasins Champion, il s'agit d'une fusion avec certains magasins Carrefour. Mme Thiran doit donc gérer les différences d'opinion et les tensions qui peuvent déjà exister entre les deux groupes.\newline \todo[inline]{Parler du fait que Thiran se voit elle-même comme plus proche des gens?=> Motivatrice, bonne humeur etc}

La comparaison du rôle de M. Nolf et M. Goergen est plus simple car ils partagent peu de points communs. En effet, la Police et IBA sont des entreprises bien différentes, qui ont des défis aujourd'hui qui sont bien distincts. IBA se porte bien économiquement ce qui induit qu'elle doit effectuer moins de changements mais plus s'orienter vers une stratégie de croissance alors que la police subit des réductions budgétaires et doit améliorer sa rentabilité avec moins d'effectifs. Ils ne font pas face aux mêmes défis ni au même type de personnel. La même analyse peut se faire entre M. Nolf et Mme Thiran, puisqu'elle travaille actuellement sur un renouvellement de culture organisationnelle, tandis que M. Nolf continue de bâtir une culture organisationnelle déjà existante dans laquelle les nouveaux membres doivent s'intégrer. La première doit recycler deux cultures pour en bâtir une nouvelle, là où le second peut s'appuyer sur une culture existante à améliorer et innover.


%Sorte de conclusion de la comparaison
Suite à cette comparaison, on se rend vite compte que chaque DRH a un rôle bien particulier en fonction de son entreprise, mais aussi et surtout en fonction de la conjoncture de l'entreprise. Nous pouvons encore une fois rappeler que ces DRH ont eu d'autres rôles au cours de leur carrière. 
\todo[inline]{Pas compris la dernière phrase}


\section{Rôles des DRH dans le modèle d'Ulrich}
%Comparaison avec modele ulrich
\section{Rôles des DRH dans le modèle d'Ulrich}

Nous avons à présent une vision des différents rôles que les intervenants endossent. Comme nous le savons, le modèle d'Ulrich donne quatre rôles que peut avoir un responsable de ressources humaines. Il peut donc être intéressant de replacer dans ce modèle, les différents DRH que nous avons rencontré. \\

Commençons par Alain Goergen. Il nous a décrit sa mission comme étant la restructuration complète du département des ressources humaines, et par la suite d'autres départements de la Police Fédérale. Mais il doit aussi aider les employés à s'habituer à ce changement. Au vue de ce rôle qu'il doit jouer, on voit très clairement que cela correspond à l'\og agent de changement \fg{} dans le modèle d'Ulrich. Pour rappel, ce dernier a pour but de faciliter le changement, donc faire en sorte que tout le monde accepte et s'adapte au changement.\\

%En lien avec le modèle d'Ulrich, on peut commencer par définir Alain Georgen comme un partenaire stratégique. Un partenaire stratégique doit assurer que le de déploiement des réformes mises en place permettent d'atteindre les objectifs recherchés. Cette restructuration a beaucoup de retombées dans différents domaines. A cause de la baisse de budget accordé à la Police, il a fallu redéfinir avec beaucoup d'attention les compétences à décentraliser ou non, prévoir les répercussions financières de celles-ci etc. Ceci a donc nécessité un travail en amont important.\\

%Dans la continuité de cette restructuration interne, on peut définir Alain Goergen comme un agent changement. Un agent du changement est quelqu'un qui s'assure que l'évolution de l'entreprise dans laquelle il travaille se passe correctement, que les comportements des individus soient bien adaptés à l'environnement. C'est exactement ce rôle qu'a joue M. Georgen dans le cadre de cette restructuration. 

Nous pouvons aussi le voir comme \og partenaire stratégique \fg{}. En effet, la décentralisation qu'il doit opérer, a un impacte stratégique en plus d'un impact sur les employés comme présenté ci-dessus. En reprenant les mots de Mr Goergen \og Le modus était tout ce qui peut être déconcentré doit être déconcentré mais si il n’y a pas de plus value à déconcentrer il faut le garder au niveau central \fg{} \footnote{TO DO : url video + minute}. Il parle ici de \og plus-value \fg{}, il y a donc fait une analyse stratégique sur les compétences à décentraliser. Dans ce cas, il ne lui a pas été demandé d'appliquer un plan pensé en amont, mais de le créer par lui même. Ses décisions ont donc été prises en fonction de l'objectif d'efficacité à atteindre. Il est aussi important de se rappeler que la Police fait fasse à des réductions de budget et que ses décisions doivent les prendre en compte. Cela fait donc de lui un \og partenaire stratégique \fg{} au sens d'Ulrich.\\   

\todo[inline]{This is a shaft}

Pour ce qui est de Frédéric Nolf, GRH d'IBA, on peut le classer lui aussi en tant que \textbf{agent de changement} dans le modèle d'Ulrich. Le rôle qu'il a joué durant la constitution d'une dream team pour le nouveau CEO permet de le classer assez facilement dans cette catégorie. Mais ce qui conforte ceci est la façon dont Nolf profile les DRH de nos jours:\newline

\og{} \textit{Je pense que les DRH modèles [\ldots] se profilent tous comme 
des architectes et des transformateurs d’entreprise \ldots} \fg{} \newline

Mais Nolf correspond aussi à un \textbf{champion des employés} dans le modèle d'Ulrich dans la façon dont il traite le personnel d'IBA. En effet, il prend en compte les avis des gens afin que ces derniers ne se sentent pas lésés et restent motivés. De plus, la création d'activités communes comme le sport ou bien de cellules de réflexion sur le développement durable permettent aux employés de s'épanouir au sein de la société. On a donc une conception très humaine de l'employé. On veut que chaque ingénieur et chercheur se sente bien dans la société. \newline

Enfin, on peut considérer le GRH d'IBA comme un \textbf{partenaire stratégique}. Un des soucis de Frédéric Nolf, comme indiquer plus haut, est de garder les employés concentrés dans leurs tâches. Les chercheurs sont des gens qui font facilement des choses à leur manière. De plus, une autre de ses préoccupations est d'exploiter le potentiel de chaque employé malgré l'internationalisation. Tout ceci est dans l'optique d'augmenter la productivité, l'efficacité de l'entreprise. Produire plus mais avec le même effectif. 



%Conclusion du la comparaison au modèle 

\end{document}
